\section{Method}

\subsection{Monotonic rational-quadratic transforms}
\label{sec:rq-splines}

We propose to implement the function $ g_{\bftheta_{i}} \!$ using \emph{monotonic rational-quadratic splines} as a building block, where each bin is defined by a monotonically-increasing rational-quadratic function. 
A rational-quadratic function takes the form of a quotient of two quadratic polynomials.
Rational-quadratic functions are easily differentiable, and since we consider only monotonic segments of these functions, they are also analytically invertible. 
Nevertheless, they are strictly more flexible than quadratic functions, and allow direct parameterization of the derivatives and heights at each knot.
In our implementation, we use the method of \citet{gregory:1982:rationalquadratic} to parameterize a monotonic rational-quadratic spline. 
The spline itself maps an interval $ [-B, B] $ to $ [-B, B] $. We define the transformation outside this range as the identity, resulting in linear `tails', so that the overall transformation can take unconstrained inputs.

The spline uses $K$ different rational-quadratic functions, with boundaries set by \mbox{$K\!+\!1$} coordinates $\curlybr{(x^{(k)}, y^{(k)})}_{k=0}^{K}$ known as \emph{knots}. The knots monotonically increase between $(x^{(0)}, y^{(0)}) \!=\! (-B, -B)$ and $(x\raisebox{-1pt}{\small ${}^{(K)}$}, y\raisebox{-1pt}{\small ${}^{(K)}$}) \!=\! (B, B)$. We give the spline $ K \!-\! 1 $ arbitrary positive values $ \curlybr{\delta^{(k)}}_{k=1}^{K-1} $ for the derivatives at the internal points, and set the boundary derivatives $ \delta^{(0)} = \delta^{(K)} = 1 $ to match the linear tails. 
If the derivatives are not matched in this way, the transformation is still continuous, but its derivative can have jump discontinuities at the boundary points.
This in turn makes the log-likelihood training objective discontinuous, which in our experience manifested itself in numerical issues and failed optimization.

The method constructs a monotonic, continuously-differentiable, rational-quadratic spline which passes through the knots, with the given derivatives at the knots. 
Defining $ s_{k} = \roundbr{y^{k+1} - y^{k}}/\roundbr{x^{k+1} - x^{k}} $ and $ \xi(x) = (x - x^{k}) / (x^{k+1} - x^{k}) $, the expression for the rational-quadratic $ \alpha^{(k)}(\xi) / \beta^{(k)}(\xi) $ in the $ k^{\text{th}} $ bin can be written
\begin{align}
    \frac{\alpha^{(k)}(\xi)}{\beta^{(k)}(\xi)} = y^{(k)} + \frac{(y^{(k + 1)} - y^{(k)}) \squarebr{s^{(k)} \xi^{2} + \delta^{(k)} \xi (1 - \xi)}}{s^{(k)} + \squarebr{\delta^{(k + 1)} + \delta^{(k)} - 2 s^{(k)}} \xi (1 - \xi)}.
    \label{eq:rq-transformation}
\end{align}
Since the rational-quadratic transformation acts elementwise on an input vector and is monotonic, the logarithm of the absolute value of the determinant of its Jacobian can be computed as the sum of the logarithm of the derivatives of \cref{eq:rq-transformation} with respect to each of the transformed $ x $ values in the input vector. 
It can be shown that
\begin{align}
	\frac{\diff}{\diff x} \squarebr{ \frac{\alpha^{(k)}(\xi)}{\beta^{(k)}(\xi)} } =  \frac{\roundbr{s^{(k)}}^{2} \squarebr{ \delta^{(k + 1)} \xi^{2} + 2 s^{(k)} \xi (1 - \xi) + \delta^{(k)} (1 - \xi)^{2} }}{\squarebr{s^{(k)} + \squarebr{\delta^{(k + 1)} + \delta^{(k)} - 2 s^{(k)}} \xi (1 - \xi)}^{2}}.
	\label{eq:rq-derivative}
\end{align}
Finally, the inverse of a rational-quadratic function can be computed analytically by inverting \cref{eq:rq-transformation}, which amounts to solving for the roots of a quadratic equation.
Because the transformation is monotonic, we can always determine which of the two quadratic roots is correct, and that the solution is given by $ \xi(x) = 2c / \roundbr{- b - \sqrt{ b^{2} - 4 a c}} $, where
\begin{align}
    a &=  \roundbr{y^{(k + 1)} - y^{(k)}} \squarebr{s^{(k)} - \delta^{(k)}} + \roundbr{ y - y^{(k)} } \squarebr{\delta^{(k + 1)} + \delta^{(k)} - 2 s^{(k)}},\\
	b &= \roundbr{y^{(k + 1)} - y^{(k)}} \delta^{(k)} -  \roundbr{ y - y^{(k)} } \squarebr{\delta^{(k + 1)} + \delta^{(k)} - 2 s^{(k)}},\\
	c &= - s^{(k)} \roundbr{ y - y^{(k)} },
\end{align}
which can the be used to determine $ x $.
An instance of the rational-quadratic transform is illustrated in \cref{fig:rq-spline}, and
\cref{appendix:rq-transform-forward} gives full details of the above expressions.

\paragraph{Implementation} The practical implementation of the monotonic rational-quadratic coupling transform is as follows:
\begin{enumerate}[topsep=0pt, itemsep=1pt]
    \item A neural network NN takes $\bfx_{1:d-1}$ as input and outputs an unconstrained parameter vector $ \bftheta_{i} $ of length $ 3K - 1$ for each $i=d, \dots, D$.
    \item Vector $ \bftheta_{i} $ is partitioned as $ \bftheta_{i} =\squarebr{\bftheta_{i}^{w},\bftheta_{i}^{h}, \bftheta_{i}^{d}}$, where $\bftheta_{i}^{w}$ and $\bftheta_{i}^{h}$ have length $K$, and $\bftheta_{i}^{d}$ has length $ K - 1 $.
    \item Vectors $\bftheta_{i}^{w}$ and $\bftheta_{i}^{h}$ are each passed through a softmax and multiplied by $2B$; the outputs are interpreted as the widths and heights of the $ K $ bins, which must be positive and span the $[-B, B]$ interval. Cumulative sums of the $K$ bin widths and heights, each starting at $-B$, yield the $ K\!+\!1 $ knots $\curlybr{(x^{(k)}, y^{(k)})}_{k=0}^K$. 
    \item The vector $\bftheta_{i}^{d}$ is passed through a softplus function and is interpreted as the values of the derivatives $ \curlybr{\delta^{(k)}}_{k=1}^{K-1} $ at the internal knots.
\end{enumerate}
Evaluating a rational-quadratic spline transform at location $ x $ requires finding the bin in which $ x $ lies, which can be done efficiently with binary search, since the bins are sorted. 
The Jacobian determinant can be computed in closed-form as a product of quotient derivatives, while the inverse requires solving a quadratic equation whose coefficients depend on the value to invert;  we provide details of these procedures in \cref{appendix-rq-transform-derivative} and \cref{appendix-rq-transform-inverse}.
Unlike the additive and affine transformations, which have limited flexibility, a differentiable monotonic spline with sufficiently many bins can approximate any differentiable monotonic function on the specified interval $ [-B, B] $\footnote{By definition
of the derivative, a differentiable monotonic function is locally linear everywhere, and can thus be approximated by a piecewise linear function arbitrarily well given sufficiently many bins. For a fixed and finite number of bins such universality does not hold, but this limit argument is similar in spirit to the universality proof of \citet{Huang:2018:naf}, and the universal approximation capabilities of neural networks in general.}, yet has a closed-form, tractable Jacobian determinant, and can be inverted analytically.
Finally, our parameterization is fully-differentiable, which allows for training by gradient methods.

The above formulation can also easily be adapted for autoregressive transforms; each $ \bftheta_{i} $ can be computed as a function of $ \bfx_{1:i-1} $ using an autoregressive neural network, and then all elements of $\bfx$ can be transformed at once. 
Inspired by this, we also introduce a set of splines for our coupling layers which act elementwise on $ \bfx_{1:d-1} $ (the typically non-transformed variables), and whose parameters are optimized directly by stochastic gradient descent. 
This means that our coupling layer transforms
\newcommand{\fw}{0.31\linewidth}
\begin{wrapfigure}{rt!}{0.5\linewidth}
	\centering
	\setlength\tabcolsep{1pt} 
	\vspace{-0em}
	\begin{tabular}{ccc}
		Training data & Flow density & Flow samples \\
		\includegraphics[width=\fw]{figures/synthetic/diamond/diamond-data} &
		\includegraphics[width=\fw]{figures/synthetic/diamond/diamond-density} &
		\includegraphics[width=\fw]{figures/synthetic/diamond/diamond-samples} 
		\\
		\includegraphics[width=\fw]{figures/synthetic/shannon/shannon-data} &
		\includegraphics[width=\fw]{figures/synthetic/shannon/shannon-density} &
		\includegraphics[width=\fw]{figures/synthetic/shannon/shannon-samples}
	\end{tabular}
	\caption{Qualitative results for two-dimensional synthetic datasets using RQ-NSF with two coupling layers.
	}
	\vspace{-2em}
	\vspace{-4em}
	\label{fig:plane-fits}
\end{wrapfigure}
all elements of $\bfx$ at once as follows: 
\begin{align}
    \bftheta_{1:d-1} &= \text{Trainable parameters} \\ 
    \bftheta_{d:D} &= \text{NN}(\bfx_{1:d-1})\label{eq:coupling_nn} \\ 
    y_i &= g_{\bftheta_i}(x_i) \quad\text{for } i=1, \dots, D.
\end{align}
\Cref{fig:plane-fits} demonstrates the flexibility of our rational-quadratic coupling transform on synthetic two-dimensional datasets. 
Using just two coupling layers, each with $ K \!=\! 128 $ bins, the monotonic rational-quadratic spline transforms have no issue fitting complex, discontinuous densities with potentially hundreds of modes. 
In contrast, a coupling layer with affine transformations has significant difficulty with these tasks (see \cref{appendix-2d-datasets-affine}).


\subsection{Neural spline flows}

The monotonic rational-quadratic spline transforms described in the previous section act as drop-in replacements for affine or additive transformations in both coupling and autoregressive transforms. 
When combined with alternating invertible linear transformations, we refer to the resulting class of normalizing flows as \emph{rational-quadratic neural spline flows} (RQ-NSF), which may feature coupling layers, RQ-NSF (C), or autoregressive layers, RQ-NSF (AR)\@. 
RQ-NSF (C) corresponds to Glow~\citep{Kingma:2018:glow} with affine or additive transformations replaced with monotonic rational-quadratic transforms, where Glow itself is exactly RealNVP with permutations replaced by invertible linear transformations.
RQ-NSF (AR) corresponds to either IAF or MAF, depending on whether the flow parameterizes $ \bff $ or $ \bff^{-1} $,
again with affine transformations replaced by monotonic rational-quadratic transforms, and also with permutations replaced with invertible linear layers.
Overall, RQ-NSF resembles a traditional feed-forward neural network architecture, alternating between linear transformations and elementwise non-linearities, while retaining an exact, analytic inverse. 
In the case of RQ-NSF (C), the inverse is available in a single neural-network pass.
